\chapter{Empirical Facts of Order Books}\label{ch:mx:empirical-facts-of-order-books}

In the second quarter of 2026, among the orders on the books of the largest US stocks that were not executed, 23\% were cancelled within a
millisecond of their arrival, 54\% within a tenth of a second and 95\% within a minute (the SEC's MIDAS statistics). Nobody designed that
number, and nobody designed the others this chapter collects: that displayed depth peaks behind the best price, that new orders are placed
at distances from the best that follow a power law, that the spread moves with the volatility of each trade. They are the order book's
stylised facts (One Quant Book 7, chapter~5, for those of returns): regularities that hold across stocks, venues and years, that any model of
a book should reproduce, and that the simulated market this book uses reproduces only in part. The chapter measures both and says which facts
the simulator fails.

\section{Depth profiles}

The average book of a liquid stock is not deepest at the best price. Bouchaud, Mézard and Potters (2002) found in three Paris Bourse stocks a
humped average profile: depth rises for a few ticks behind the best and then decays, a shape a \omterm{def:mx:the-limit-order-book:zi}{zero-intelligence model} reproduces. The
simulated hour of chapter~1 has the same shape (\cref{fig:mx:the-limit-order-book:fill}, right): 1\,366 shares at the best on average, 1\,508 one
tick behind. The reason is visible in the flows. In the simulated day (6.5 hours, 715\,054 messages), cancellations per displayed share per second
are 0.090 at the best, 0.057 one tick behind and 0.044 two ticks behind: the best level is where orders are stale first and executed first, so it
holds less than the levels a \omterm{def:mx:the-limit-order-book:orders}{market order} rarely reaches.

\section{Where new orders go}

\begin{definition}[Relative limit price]\label{def:mx:empirical-facts-of-order-books:relative}
The \emph{relative limit price}\index{relative limit price} of a new \omterm{def:mx:the-limit-order-book:orders}{limit order} is its distance, in ticks, from the best price on its own side when
it arrives: positive behind the best, zero at it, negative when it improves it.
\end{definition}

Zovko and Farmer (2002) merged seven million London Stock Exchange orders on 50 stocks and found that the cumulative distribution of \omterm{def:mx:empirical-facts-of-order-books:relative}{relative limit
prices} decays roughly as a power law with exponent about 1.5, from a few ticks to about two thousand: most orders arrive near the best, but a few
are placed very far away, and their share falls slowly. Bouchaud, Mézard and Potters found a power law too. The simulator does not have one: Book~7's
\texttt{firm.tape} places \omterm{def:mx:the-limit-order-book:orders}{limit orders} on its ten best levels with rates that decay geometrically. In the simulated day 0.5\% of new orders improve the
best, 18.6\% join it, 40.1\% go four ticks or more behind and none goes more than nine. A strategy whose edge depends on deep liquidity (a large
metaorder walking the book, chapter~11) would find the simulated book too thin far from the best.

\section{Spreads and the volatility of each trade}

\begin{definition}[Volatility per trade]\label{def:mx:empirical-facts-of-order-books:vol}
The \emph{volatility per trade}\index{volatility per trade} $\sigma_1$ of an instrument is the standard deviation of the change in the mid price between
one trade and the next.
\end{definition}

Wyart, Bouchaud, Kockelkoren, Potters and Vettorazzo (2008) showed that if neither market making nor market taking earns more than a marginal profit,
the spread must be proportional to the instantaneous impact of a trade; they found the spread and the \omterm{def:mx:empirical-facts-of-order-books:vol}{volatility per trade} strongly correlated
across stocks, with $R^2$ above 0.9, and read it as evidence that adverse selection sets the spread and that most volatility comes from the impact
of trades (chapters 4, 5 and 11 take the three halves of that sentence in turn). \Cref{fig:mx:empirical-facts-of-order-books:spread} tests the relation
in two simulated markets. In chapter~1's zero-intelligence market, whose spread is free to vary, eight configurations of the rates line up with
$s \approx 0.96 + 1.64\,\sigma_1$ and $R^2 = 0.989$. In ten configurations of \texttt{firm.tape} the relation all but disappears ($R^2=0.48$, slope 0.19):
the simulated stock is a large-tick stock whose spread is one tick 96.6\% of the time, and the spread has no room to respond. Chapter~6 names the two
regimes; the relation is a small-tick fact.

\begin{omfigure}
\begin{tikzpicture}
\begin{axis}[width=9cm,height=5.8cm,xlabel={\small volatility per trade $\sigma_1$ (ticks)},ylabel={\small mean spread (ticks)},
  tick label style={font=\footnotesize},xmin=0,xmax=5.5,ymin=0,ymax=10,grid=major,grid style={gray!25},
  legend columns=3,legend style={font=\footnotesize,at={(0.5,-0.3)},anchor=north,draw=none,fill=none,
  /tikz/every even column/.append style={column sep=8pt}},table/col sep=comma]
\addplot[only marks,mark=o,omDef,thick] table[x=vol,y=spread,restrict expr to domain={\thisrow{kind}}{0:0}]{figdata/microstructure/03-empirical-facts-of-order-books/spreadvol.csv};
\addplot[only marks,mark=square*,omThm] table[x=vol,y=spread,restrict expr to domain={\thisrow{kind}}{1:1}]{figdata/microstructure/03-empirical-facts-of-order-books/spreadvol.csv};
\addplot[omDef,thick,dashed] table[x=vol,y=spread]{figdata/microstructure/03-empirical-facts-of-order-books/spreadvol_fit.csv};
\legend{zero-intelligence market,\texttt{firm.tape},fit $0.96+1.64\,\sigma_1$}
\end{axis}
\end{tikzpicture}

\omcaption{Mean spread against \omterm{def:mx:empirical-facts-of-order-books:vol}{volatility per trade}: eight configurations of chapter~1's zero-intelligence market (100\,000 events each) and ten of
\texttt{firm.tape} (20 simulated minutes each, informed intensity and efficient-price jump rate varied). The zero-intelligence spreads follow the
line; \texttt{firm.tape}'s stay near one tick. Data: \texttt{mx\_facts.spread\_vs\_vol\_zi}, \texttt{mx\_facts.spread\_vs\_vol}.}\label{fig:mx:empirical-facts-of-order-books:spread}
\end{omfigure}

\section{Intraday shapes}

Volume, messages and volatility follow the clock (One Quant Book 7, chapter~5): high at the open, lower at midday, higher again into the close, while
spreads in real books are widest in the first minutes and narrow through the morning. \texttt{firm.tape} plants only the first half: its rates are
multiplied by a U-shaped profile. One simulated day is too noisy to show it, because a persistent random activity level dominates any hour; averaged over
eight compressed sessions, volume falls from 92\,200 shares in the first thirteenth of the session to 24\,700 at the lowest and returns to 55\,600 in the
last (\cref{fig:mx:empirical-facts-of-order-books:intraday}). Its spread stays at one tick at every hour: the simulator has no opening spread.

\begin{omfigure}
\begin{tikzpicture}
\begin{axis}[width=8.8cm,height=5cm,xlabel={\small thirteenth of the session},ylabel={\small thousands},tick label style={font=\footnotesize},
  xtick={1,2,3,4,5,6,7,8,9,10,11,12,13},ymin=0,ymax=100,grid=major,grid style={gray!25},ybar,bar width=6pt,
  legend columns=2,legend style={font=\footnotesize,at={(0.5,-0.32)},anchor=north,draw=none,fill=none,
  /tikz/every even column/.append style={column sep=8pt}},table/col sep=comma]
\addplot[fill=omDef!60,draw=omDef] table[x=bin,y=volume]{figdata/microstructure/03-empirical-facts-of-order-books/intraday.csv};
\addplot[fill=omThm!60,draw=omThm] table[x=bin,y=messages]{figdata/microstructure/03-empirical-facts-of-order-books/intraday.csv};
\legend{shares traded,messages}
\end{axis}
\end{tikzpicture}

\omcaption{The simulated intraday profile: shares traded and messages in each thirteenth of the session, averaged over eight sessions of \texttt{firm.tape}
with a U-shaped activity profile. Data: \texttt{mx\_facts.intraday\_mean}.}\label{fig:mx:empirical-facts-of-order-books:intraday}
\end{omfigure}

\section{Order lifetimes, fleeting orders and resilience}

\begin{definition}[Order lifetime, fleeting order]\label{def:mx:empirical-facts-of-order-books:life}
An order's \emph{order lifetime}\index{order lifetime} is the time from its arrival to its full cancellation or execution. A \emph{fleeting
order}\index{fleeting order} is one cancelled in full within a short time of its arrival (two seconds is a common threshold).
\end{definition}

Hasbrouck and Saar (2009) studied such orders on a Nasdaq trading platform and argued that they behave more like demands for liquidity than supplies
of it: traders post and withdraw to search for hidden counterparties or to chase a moving price. The SEC's MIDAS system publishes lifetimes for every US
stock: the time from an order's arrival to a cancellation or execution, with the clock restarting after each partial event, estimated as a lifetable
in which executions censor cancellations. \Cref{fig:mx:empirical-facts-of-order-books:life} compares the cancellation curves with the simulator's. In
large stocks half the cancellations come within a tenth of a second, in small stocks within about a second; in \texttt{firm.tape} 6.8\% come within a
second and 40\% within ten. The simulator's orders live a hundred times longer than real ones: its participants are slow, and the chapters that depend on
the race between cancellations and executions (latency, chapter~8; queue positions, chapter~6) must say which market they are measured in.

\begin{omfigure}
\begin{tikzpicture}
\begin{axis}[width=9cm,height=5.4cm,xlabel={\small time since arrival or last partial event},ylabel={\small share cancelled},
  tick label style={font=\footnotesize},xtick={0,1,2,3,4,5},xticklabels={1 ms,10 ms,100 ms,1 s,10 s,1 min},ymin=0,ymax=1,
  grid=major,grid style={gray!25},legend columns=2,legend style={font=\footnotesize,at={(0.5,-0.3)},anchor=north,draw=none,fill=none,
  /tikz/every even column/.append style={column sep=8pt}},table/col sep=comma]
\addplot[omDef,very thick,mark=*] table[x=k,y=large]{figdata/microstructure/03-empirical-facts-of-order-books/lifetimes.csv};
\addplot[omProp,thick,mark=triangle*] table[x=k,y=mid]{figdata/microstructure/03-empirical-facts-of-order-books/lifetimes.csv};
\addplot[omThm,thick,mark=square*] table[x=k,y=small]{figdata/microstructure/03-empirical-facts-of-order-books/lifetimes.csv};
\addplot[black,thick,dashed,mark=o] table[x=k,y=sim]{figdata/microstructure/03-empirical-facts-of-order-books/lifetimes.csv};
\legend{large US stocks,mid-cap US stocks,small US stocks,\texttt{firm.tape} (simulated day)}
\end{axis}
\end{tikzpicture}

\omcaption{The share of orders cancelled within a given time, executions treated as censoring (lifetable estimates): US stocks by capitalisation group in
2026 Q2 (SEC MIDAS, Hazards and Survivors) and the simulated day. Data: \texttt{data/microstructure/midas\_lifetimes.csv}, \texttt{mx\_facts.lifetable}.}\label{fig:mx:empirical-facts-of-order-books:life}
\end{omfigure}

\begin{definition}[Book resilience]\label{def:mx:empirical-facts-of-order-books:resilience}
The \emph{book resilience}\index{book resilience} of a \omterm{def:mx:the-limit-order-book:lob}{limit order book} is the speed at which its spread and depth return to their usual state after a
trade or a cancellation has depleted them.
\end{definition}

Large (2007) measured it on an electronic book as the response of new orders to liquidity shocks. In the simulated day, 750 executions emptied the best
level and widened the spread to two ticks or more; the median time until the spread was one tick again was 0.19 seconds, and 92.4\% recovered within a
second. Resilience is what separates a temporary price concession from a permanent one, and chapter~12's impact kernels are resilience measured in
prices.

The facts also move with market structure, and the MIDAS statistics show how much in a decade.

\begin{dated}{2026-06}{US stock order flow in the MIDAS statistics}\label{dat:mx:empirical-facts-of-order-books:midas}
Daily means over all US equity exchanges for stocks (SEC MIDAS, Summary Metrics by Exchange): the cancel-to-trade ratio was 19.8 in 2012 and 16.6 in the
first half of 2026; the hidden rate (the share of trades executed against \omterm{def:mx:order-types-and-their-uses:hidden}{hidden orders}) rose from 10.5\% to 32.2\%; the odd-lot rate (the share of trades
smaller than a round lot) from 19.2\% to 68.4\%. In 2026 Q2, among orders in the largest stocks, 54\% of the cancellations came within 100 milliseconds of
arrival and 95\% within a minute.
\end{dated}

\begin{omfigure}
\begin{tikzpicture}
\begin{axis}[name=rates,width=7cm,height=5cm,title={\small share of trades (\%)},tick label style={font=\footnotesize},
  xtick={2012,2014,2016,2018,2020,2022,2024,2026},x tick label style={/pgf/number format/1000 sep={},rotate=45,anchor=east},
  ymin=0,ymax=75,grid=major,grid style={gray!25},legend columns=2,legend style={font=\footnotesize,at={(1.1,-0.35)},anchor=north,draw=none,
  fill=none,/tikz/every even column/.append style={column sep=8pt}},table/col sep=comma]
\addplot[omDef,very thick,mark=*] table[x=year,y=oddlot]{figdata/microstructure/03-empirical-facts-of-order-books/yearly.csv};
\addplot[omThm,thick,mark=square*] table[x=year,y=hidden]{figdata/microstructure/03-empirical-facts-of-order-books/yearly.csv};
\legend{odd-lot rate,hidden rate}
\end{axis}
\begin{axis}[at={(rates.east)},anchor=west,xshift=1.2cm,width=7cm,height=5cm,title={\small cancel-to-trade ratio},tick label style={font=\footnotesize},
  xtick={2012,2014,2016,2018,2020,2022,2024,2026},x tick label style={/pgf/number format/1000 sep={},rotate=45,anchor=east},
  ymin=0,ymax=25,grid=major,grid style={gray!25},table/col sep=comma]
\addplot[black,thick,mark=o] table[x=year,y=ctt]{figdata/microstructure/03-empirical-facts-of-order-books/yearly.csv};
\end{axis}
\end{tikzpicture}

\omcaption{US stocks, yearly means of daily statistics over all exchanges, 2012 to June 2026: the share of trades in odd lots and against \omterm{def:mx:order-types-and-their-uses:hidden}{hidden orders}
(left), and cancellations per trade (right). Data: SEC MIDAS, Summary Metrics by Exchange (\texttt{data/microstructure/midas\_yearly.csv}).}\label{fig:mx:empirical-facts-of-order-books:yearly}
\end{omfigure}

\section{Is the simulator a market?}

The table sets the chapter's facts beside the simulator's numbers. The simulated day's 15.5 cancellations per trade happen to match the MIDAS ratio of
recent years; most of the rest do not.

\begin{center}
\small
\begin{tabular}{@{}p{3.6cm}p{4.6cm}p{3.8cm}l@{}}
\toprule
fact & real books & \texttt{firm.tape} & holds? \\
\midrule
depth profile & humped (Paris Bourse) & peak one tick behind the best & yes \\
\omterm{def:mx:empirical-facts-of-order-books:relative}{relative limit prices} & power law, exponent $\approx1.5$, to 2\,000 ticks (London) & none beyond nine ticks & no \\
spread and \omterm{def:mx:empirical-facts-of-order-books:vol}{volatility per trade} & $R^2>0.9$ across stocks & $R^2=0.48$ (large tick; the zero-intelligence market: 0.989) & no \\
intraday volume & U-shaped & U-shaped on average (planted) & yes \\
opening spread & widest early in the day & one tick all day & no \\
cancellation speed & 54\% within 0.1 s (large US stocks) & 0.8\% within 0.1 s & no \\
cancel-to-trade ratio & 16.6 (US stocks, 2026) & 15.5 & yes \\
hidden liquidity & 32\% of trades (2026) & none & no \\
resilience after a depleting trade & fast (sub-second in liquid stocks) & median 0.19 s & yes \\
\bottomrule
\end{tabular}
\end{center}

A simulator used to measure queue positions, latency or hidden liquidity must fix its failing rows first; chapter~27 builds an agent-based market
calibrated to these facts with \texttt{firm.lobstats}.

\section{Tutorial: measure the facts}\label{tut:mx:empirical-facts-of-order-books:tut}

\begin{tutorial}
\textbf{Goal.} Measure the stylised facts on the simulated day and set them beside the SEC's statistics. \textbf{End state:} the table and
\cref{fig:mx:empirical-facts-of-order-books:spread,fig:mx:empirical-facts-of-order-books:life}.
\begin{tutsteps}
\item \textbf{The public data.} \texttt{mx\_fetch\_midas.py} downloads the MIDAS hazard and exchange files for 2026 Q2 and writes the two derived tables to
\texttt{data/microstructure/} (the committed copies are enough for the rest).
\item \textbf{Lifetimes.} The lifetable with executions as censoring, as MIDAS computes it:
\omcode{code/microstructure/03-empirical-facts-of-order-books/python/mx_facts.py}{121}{141}{The lifetable estimate of time to cancellation.}\label{lst:mx:empirical-facts-of-order-books:life}
\item \textbf{Flows and shapes.} \texttt{relative\_prices}, \texttt{cancel\_rates}, \texttt{spread\_profile}, \texttt{intraday\_mean}, \texttt{resilience}
and \texttt{fleeting\_share} on \texttt{day()}.
\item \textbf{The relation.} \texttt{spread\_vs\_vol\_zi()} and \texttt{spread\_vs\_vol()}; draw with \texttt{fig\_facts.py}.
\end{tutsteps}
\textbf{What to change next.} Give \texttt{firm.tape} a tick of half a cent (halve every price) and rerun the spread relation; replace its geometric
placement by a power law and measure the depth profile again.
\end{tutorial}

\section{Build: the stylised-fact toolkit}\label{bld:mx:empirical-facts-of-order-books:lobstats}

\begin{build}
\bfield{Purpose} One report of an order book's stylised facts, run on any market-by-order stream: the simulator's, a recorded feed, a generated one
(Book~12's synthetic data); the target chapter~27 calibrates against.
\bfield{Interface} \texttt{firm\_lobstats.report(tape, levels, points)} returning the depth profile, the relative-price distribution,
cancellation rates by distance, the spread distribution, the \omterm{def:mx:empirical-facts-of-order-books:vol}{volatility per trade}, lifetable curves (cancellation and execution), the fleeting share,
resilience and the cancel-to-trade ratio; \texttt{kaplan\_meier(durations, observed, points)}; \texttt{compare(report, targets)} checking each fact
against a range.
\bfield{Rules} Inputs in \texttt{firm.tape}'s message format (the simulator's \texttt{res.tape()} produces it); lifetimes restart after partial events and
executions censor cancellations, as in MIDAS; time-weighted averages for the spread and depth.
\bfield{Acceptance tests} \texttt{code/firm/lobstats/tests/}: the product-limit estimate by hand; a simulated session's humped depth, one-tick spread,
monotone lifetables and fast resilience; \texttt{compare} flagging the simulator's slow cancellations.
\bfield{Stretch} A power-law tail fit with a Hill estimator (One Quant Book 4, chapter~15) and its confidence interval; the same report on a recorded
day from \texttt{make\_recorded\_day.py}.
\end{build}

\begin{omsources}
\item J.-P.~Bouchaud, M.~Mézard and M.~Potters, ``Statistical properties of stock order books: empirical results and models'', \emph{Quantitative
Finance} 2, 2002.
\item I.~Zovko and J.~D.~Farmer, ``The power of patience: a behavioural regularity in limit-order placement'', \emph{Quantitative Finance} 2, 2002.
\item M.~Wyart, J.-P.~Bouchaud, J.~Kockelkoren, M.~Potters and M.~Vettorazzo, ``Relation between bid-ask spread, impact and volatility in order-driven
markets'', \emph{Quantitative Finance} 8(1), 2008.
\item J.~Hasbrouck and G.~Saar, ``Technology and liquidity provision: the blurring of traditional definitions'', \emph{Journal of Financial Markets}
12(2), 2009.
\item A.~J.~Large, ``Measuring the resiliency of an electronic limit order book'', \emph{Journal of Financial Markets} 10(1), 2007.
\item B.~Biais, P.~Hillion and C.~Spatt, ``An empirical analysis of the limit order book and the order flow in the Paris Bourse'', \emph{Journal of
Finance} 50(5), 1995.
\item SEC, MIDAS market structure data: Hazards and Survivors by Time Period; Summary Metrics by Exchange; quote-life methodology.
\end{omsources}

\section{Exercises}

\begin{exercise}[$\star$]\label{exo:mx:empirical-facts-of-order-books:1}
A buy \omterm{def:mx:the-limit-order-book:orders}{limit order} arrives at 99.97 when the best bid is 100.00 and the best ask 100.01, with a one-cent tick. What is its \omterm{def:mx:empirical-facts-of-order-books:relative}{relative limit price}? And
for a buy at 100.00?
\end{exercise}

\begin{exercise}[$\star$]\label{exo:mx:empirical-facts-of-order-books:2}
In a small-tick market the \omterm{def:mx:empirical-facts-of-order-books:vol}{volatility per trade} is 0.9 ticks. What spread does the zero-intelligence fit of the chapter predict?
\end{exercise}

\begin{exercise}[$\star$]\label{exo:mx:empirical-facts-of-order-books:3}
Why is the hidden rate of \cref{fig:mx:empirical-facts-of-order-books:yearly} a share of trades and not of orders, and what would the share of orders
measure instead?
\end{exercise}

\begin{exercise}[$\star\star$]\label{exo:mx:empirical-facts-of-order-books:4}
A trader sees 1\,000 cancellations in a large stock. According to the 2026 Q2 lifetable, how many would you expect to have come within 100 milliseconds
of the orders' arrival? What does ``executions treated as censoring'' change in that reading?
\end{exercise}

\begin{exercise}[$\star\star$]\label{exo:mx:empirical-facts-of-order-books:5}
The odd-lot rate rose from 19\% of trades in 2012 to 68\% in 2026. Give two structural reasons, and one consequence for a feed handler that ignores odd
lots.
\end{exercise}

\begin{exercise}[$\star\star$]\label{exo:mx:empirical-facts-of-order-books:6}
Why does a large-tick stock break the spread--volatility relation, and what replaces the spread as the thing that adjusts?
\end{exercise}

\begin{exercise}[$\star\star\star$]\label{exo:mx:empirical-facts-of-order-books:7}
\emph{Coding.} Compute the lifetable of the time to \emph{execution} on the simulated day (\texttt{lifetable} with executions as the event, cancellations
censoring) at 1, 10 and 60 seconds, and compare it with the MIDAS figures for large stocks: 3.4\%, 8.7\% and 15.0\%. What does the comparison say?
\end{exercise}

\begin{exercise}[$\star\star\star$]\label{exo:mx:empirical-facts-of-order-books:8}
\emph{Find the flaw.} ``Our passive strategy fills 12\% of its orders in the simulator and loses nothing to faster traders, so it will do the same live.''
\end{exercise}

\section{Problem: Is the Simulator a Market?}

\begin{problem}[{Weekend problem --- is the simulator a market?}]\label{pb:mx:empirical-facts-of-order-books:1}
Before trusting \texttt{firm.tape} for a study of passive execution, a researcher checks it against the published facts and the SEC's statistics.

\medskip\noindent\textbf{Part I --- The book.}
\begin{enumerate}
\item Where does the simulated depth profile peak, and what do real books show?
\item What are the simulated cancellation rates per displayed share at the best and two ticks behind, and why do they differ?
\item What share of new orders improves the best, joins it, or goes four ticks or more behind?
\item What do real \omterm{def:mx:empirical-facts-of-order-books:relative}{relative limit prices} look like, and what does the simulator do instead?
\end{enumerate}

\medskip\noindent\textbf{Part II --- Spreads and clocks.}
\begin{enumerate}[resume]
\item How often is the simulated spread one tick, and what is its mean?
\item What relation between spread and \omterm{def:mx:empirical-facts-of-order-books:vol}{volatility per trade} does the zero-intelligence market give, with what $R^2$?
\item What does \texttt{firm.tape} give, and why?
\item What intraday volume profile does the simulator show, and how many sessions were needed to see it?
\end{enumerate}

\medskip\noindent\textbf{Part III --- Lifetimes.}
\begin{enumerate}[resume]
\item How does MIDAS compute its lifetime curves?
\item What share of cancellations in large, and in small, US stocks comes within 100 milliseconds and within a second?
\item What are the simulator's shares within 0.1 and 1 second, and within 10 seconds?
\item What share of the simulated day's orders is fleeting (cancelled within two seconds)?
\end{enumerate}

\medskip\noindent\textbf{Part IV --- The verdict.}
\begin{enumerate}[resume]
\item How fast does the simulated book recover from a trade that empties its best level?
\item How do the simulated day's cancellations per trade compare with MIDAS?
\item State the \emph{named result}: the table of stylised facts, simulated against published, and which the simulator fails.
\item Which failures matter for a study of queue position and fill rates, and in which direction do they bias it?
\item Which failure matters for a study of large orders?
\item What changed in US order flow from 2012 to 2026, in MIDAS's figures?
\item What would you change in the simulator first?
\item In one sentence: what is a stylised fact for?
\end{enumerate}
\end{problem}

\section{Interview questions}

\begin{interviewq}[$\star$ \iqroles{researcher}]\label{iq:mx:empirical-facts-of-order-books:1}
Why is the average order book deeper one or two ticks behind the best price than at it?
\end{interviewq}

\begin{interviewq}[$\star\star$ \iqroles{researcher, trader}]\label{iq:mx:empirical-facts-of-order-books:2}
Across stocks, spreads are strongly correlated with \omterm{def:mx:empirical-facts-of-order-books:vol}{volatility per trade}. What does that tell you about what sets the spread?
\end{interviewq}

\begin{interviewq}[$\star\star$ \iqroles{researcher}]\label{iq:mx:empirical-facts-of-order-books:3}
Most orders are cancelled within a second. How would you estimate the distribution of \omterm{def:mx:empirical-facts-of-order-books:life}{order lifetimes} when some orders are executed first?
\end{interviewq}

\begin{interviewq}[$\star\star$ \iqroles{developer, researcher}]\label{iq:mx:empirical-facts-of-order-books:4}
You are given a simulator of order flow. How would you decide whether it is realistic enough to backtest a market-making strategy?
\end{interviewq}

\begin{interviewq}[$\star\star$ \iqroles{trader}]\label{iq:mx:empirical-facts-of-order-books:5}
Why are spreads wider at the open than at noon?
\end{interviewq}

\begin{interviewq}[$\star\star\star$ \iqroles{researcher}]\label{iq:mx:empirical-facts-of-order-books:6}
What would you expect a \omterm{def:mx:empirical-facts-of-order-books:life}{fleeting order} to reveal about its sender, and how would you test it?
\end{interviewq}
